% Artifact Appendix for:
%   ST-Flow: A Hardware Compiler for Spatial-Temporal Dataflow Acceleration (PACT '26)
%
% Follows the cTuning/ACM template:
%   https://github.com/ctuning/artifact-evaluation/blob/master/docs/template/ae.tex
%
% To include in the camera-ready, \input this file after \bibliography, or paste the
% body below into main.tex. It expects \usepackage{hyperref} (main.tex already loads it).
%
% This is the short-form derivative of ARTIFACT_EVALUATION.md and docs/ae/checklist.md;
% if any of the three ever disagree, ARTIFACT_EVALUATION.md is authoritative.

\appendix
\section{Artifact Appendix}

\subsection{Abstract}

The artifact contains the ST-Flow compiler---a Python driver orchestrating an MLIR-based
C++ backend that lowers PyTorch models and affine programs to FPGA streaming dataflow
accelerators---together with the full archived record of the paper's 425 experiments
(5.6\,GB, 41{,}874 files: generated sources, solver traces, HLS/implementation reports and
on-board execution logs). Reproducing every experiment from source needs days of licensed
Vivado time and a Xilinx Alveo U55C, so a single command instead replays the analysis
pipeline over the archives and regenerates every figure and table in the paper in about
eleven seconds: the consolidated CSVs and \LaTeX{} tables come back byte-identical, and the
figure PDFs differ only in their embedded creation timestamp. The artifact is organized in
four tiers, so a reviewer with no FPGA and no HLS licence can still validate the paper's
numbers end to end, and can additionally build the compiler and run its 66-test suite with
only Vitis HLS. The artifact is available for download from Zenodo at \url{https://doi.org/10.5281/zenodo.21881993}. We are applying for the \emph{Available}, \emph{Functional} and
\emph{Reproduced} badges.

\subsection{Artifact check-list (meta-information)}

{\small
\begin{itemize}
  \item {\bf Algorithm: } Spatial-temporal hardware mapping via structural/temporal
    kernel clustering; non-divisor loop tiling; design space exploration by MINLP,
    simulated annealing and genetic algorithm over a closed-form analytical latency and
    resource model.
  \item {\bf Program: } No standard suite. Six DNN workloads (MLP-Mixer, RWKV
    channel-mix, TabM, SwiGLU, DIN, MNIST MLP), eight synthetic dataflow topologies,
    seven Polybench kernels, and GPT-2 Tiny as a case study --- all PyTorch modules
    under \texttt{test/}, none need downloading.
  \item {\bf Compilation: } LLVM/MLIR~20 with Python bindings (prebuilt release
    provided); CMake~$\geq$~3.20 and Ninja; Python~3.12. Vitis HLS 2024.2 compiles the
    generated accelerators.
  \item {\bf Transformations: } The ST-Flow compiler itself: an MLIR dialect with DFG
    solvers, transform passes and HLS/TAPA emission. Public, included, at the
    evaluated revision.
  \item {\bf Model / Data set: } Weights are seeded-random
    (\texttt{torch.manual\_seed(42)}); inputs are synthetic tensors from the same
    fixed seed. Nothing is downloaded.
  \item {\bf Run-time environment: } Linux x86\_64, Python~3.12 with pinned
    \texttt{torch}/\texttt{torch\_mlir} wheels. Vitis HLS/Vivado 2024.2 and XRT for
    synthesis and on-board runs; AMPL (Gecode/Gurobi) for the MINLP solver. The
    archived-data reproduction path needs Python only.
  \item {\bf Hardware: } Xilinx Alveo U55C at 300\,MHz for bitstream generation and
    on-board execution (Tier~4 only). Tiers~1--2 need only a Linux workstation.
  \item {\bf Run-time state: } The Tier~2 path is insensitive to machine load. The
    Tier~1 test suite runs HLS C-simulation in parallel and is load-sensitive; keep the
    worker count at or below the core count.
  \item {\bf Metrics: } Cycle counts (analytical prediction, Vitis min/max estimate,
    RTL simulation, on-board); MAPE and Pearson correlation against RTL; speedup
    relative to divisor-only tiling; DSP/BRAM/URAM utilization; area-delay product;
    achieved clock frequency; DSE solution cost, wall-clock search time and
    coefficient of variation.
  \item {\bf Output: } Consolidated CSVs (\texttt{main.csv}, \texttt{gpt2.csv},
    \texttt{codo.csv}, \texttt{gpt2\_dse\_traces.csv}), seven figure PDFs and four
    \LaTeX{} tables, corresponding to Figures~3--6 and Tables~4--7 of the paper.
  \item {\bf Experiments: } 171 main-sweep designs (three tiling strategies
    $\times$ 14 workloads at three hardware-mapping strategies, plus Polybench and an
    extra DSP-budget sweep for the CODO comparison), 240 GPT-2 Tiny DSE runs
    (4 solvers $\times$ 20 repeats $\times$ 3 mappings), and 14 augmented CODO
    baseline designs (9 more failed and are archived separately). Driven by
    \texttt{PACT26\_Artifacts/experiments.sh}; checked by \texttt{scripts/ae\_verify.sh}.
  \item {\bf How much disk space required (approximately)?: } $\sim$5.6\,GB archive
    ($\sim$10\,GB after extraction), plus $\sim$5\,GB for the LLVM/MLIR install and
    $\sim$5\,GB for the compiler build. Allow 20\,GB.
  \item {\bf How much time is needed to prepare workflow (approximately)?: } About
    10~minutes for the archived-data path (clone plus \texttt{pip install}); about
    15--30~minutes more to build the compiler using the prebuilt LLVM.
  \item {\bf How much time is needed to complete experiments (approximately)?: }
    Tier~2 (all paper figures and tables): $\sim$11~seconds. Tier~1 (66-test suite):
    $\sim$4~minutes at 11 workers on a 20-core host. Tier~3 (one design through HLS):
    minutes to $\sim$1.5~hours per design. Tier~4 (bitstream and on-board): hours per design, days
    for the full sweep.
  \item {\bf Publicly available?: } Yes, as an Open Access archival record: \url{https://doi.org/10.5281/zenodo.21881993}.
    Self-contained: source, all archived experiments, the prebuilt LLVM/MLIR toolchain
    and the pinned Python wheels.
  \item {\bf Code licenses (if publicly available)?: } Apache License 2.0.
  \item {\bf Data licenses (if publicly available)?: } Creative Commons Attribution 4.0
    for the archived experiment outputs.
  \item {\bf Workflow automation framework used?: } No third-party framework. SLURM
    for the cluster sweep (not needed to evaluate the artifact).
  \item {\bf Archived (provide DOI)?: } \url{https://doi.org/10.5281/zenodo.21881993}
\end{itemize}
}

\subsection{Description}

\subsubsection{How to access}

Download the archived record from Zenodo at \url{https://doi.org/10.5281/zenodo.21881993}. It is self-contained: the
compiler at the evaluated revision, \texttt{PACT26\_Artifacts/} (the 5.6\,GB experiment
archive), the prebuilt LLVM/MLIR~20 toolchain, and the pinned Python wheels that are no
longer obtainable from any index. Nothing is fetched over the network during evaluation;
the development repository is not public, so this record is the single point of access.
Start from \texttt{ARTIFACT\_EVALUATION.md} at the top level, which states the supported
claims and gives the tiered instructions summarized below.

\subsubsection{Hardware dependencies}

A Linux x86\_64 machine suffices for Tiers~1--2; the reproduction path is CPU-bound and
needs no accelerator. Bitstream generation and on-board validation (Tier~4) require a
Xilinx Alveo U55C, the card the paper's physical results were measured on. Reviewers
without one can stop after Tier~3: the on-board claims are supported by archived
implementation reports and execution logs.

\subsubsection{Software dependencies}

Python~3.12 (the pinned wheels are \texttt{cp312}-only) and LLVM/MLIR~20 with Python
bindings, provided prebuilt since ST-Flow does not compile against distribution
LLVM~17/18. Vitis HLS~2024.2 for C-simulation and synthesis, Vivado~2024.2 and XRT for
bitstreams and on-board runs, and AMPL for the MINLP solver (a heuristic solver removes
this dependency). Two pins matter: \texttt{pybind11$<$3} (3.x cannot compile the MLIR~20
bindings) and \texttt{matplotlib==3.10.5} (the version that produced the paper's
figures); both are recorded in \texttt{requirements.txt}. The ML front-end is pinned to a
PyTorch nightly that upstream no longer serves, so the exact wheels are distributed with
the artifact.

\subsubsection{Data sets and Models}

No external data sets or model checkpoints. All inputs are synthetic tensors from a
fixed seed and all weights are seeded-random, so every experiment runs from a clean
checkout with no downloads. The workloads are listed under Program above.

\subsection{Installation}

\begin{verbatim}
tar xzf st-flow-pact26.tar.gz && cd st-flow-pact26
python3.12 -m venv .venv && . .venv/bin/activate
pip install -r requirements.txt
source env.sh
\end{verbatim}

Sufficient to reproduce every figure and table (Tier~2). To build the compiler itself
(Tiers~1, 3, 4):

\begin{verbatim}
mkdir -p ~/opt && tar xzf llvm-stflow-ubuntu-latest.tar.gz -C ~/opt
bash scripts/install_ml_wheels.sh wheels/
source env.sh
export NUM_THREADS=6      # bounded on purpose: one job per core OOMs the build
python -m pip install -v -e .
\end{verbatim}

\texttt{README.md} documents the full procedure and troubleshooting.

\subsection{Experiment workflow}

\texttt{PACT26\_Artifacts/experiments.sh} is the script that produced the paper's data:
it consolidates per-design reports into four CSVs, then drives \texttt{evplot} and
\texttt{evtable} to emit each figure and table. It reads only archived reports and logs,
so it needs no FPGA, HLS licence, solver or scheduler. \texttt{scripts/ae\_verify.sh}
wraps it: regenerate into a scratch directory, leave the shipped reference untouched,
diff the two.

Generating new designs rather than replaying archived ones goes through \texttt{pytest}
for a single design's DSE/C-simulation/synthesis/co-simulation (Tier~3), and that
design's own \texttt{Makefile} (\texttt{make bitstream}) plus
\texttt{STFLOW\_MODE=bitstream pytest} for its bitstream and on-board run (Tier~4). For
the full design grid, \texttt{evgen} plus \texttt{evrun submit} drive the same steps as
SLURM jobs.

\subsection{Evaluation and expected results}

\begin{verbatim}
source env.sh
bash scripts/ae_verify.sh
\end{verbatim}

Expected, in about eleven seconds: \texttt{RESULT: PASS} with 15 reproduced outputs and
no mismatches --- the four CSVs and four \LaTeX{} tables byte-identical to the shipped
reference, the seven figure PDFs identical in content with only their embedded creation
timestamp differing. This validates the model-validation accuracy of
Section~5.2 (MAPE 8.24\% on DNN, 17.82\% on synthetic, against 28.46--86.38\% for Vitis
estimates), the non-divisor tiling speedups of Section~5.3
($3.80\times$/$1.44\times$/$1.12\times$ for max-reuse/balanced/no-reuse), the area-delay
comparison of Section~5.4, the $23.11\times$ latency and $10.65\times$ PPA advantage over
CODO in Section~5.5, the GPT-2 DSE comparison of Section~5.6.1, and the on-board results
of Section~5.6.2. Table~7 --- the one table the paper carries inline rather than
\texttt{\textbackslash input}-ing --- is generated too and reproduces cell for cell. See
\texttt{ARTIFACT\_EVALUATION.md} for the full C1--C7 claims table.

Independently, \texttt{pytest test/} exercises the compiler itself: all 66 tests compile
PyTorch models through the full pipeline and validate each design by HLS C-simulation
against a PyTorch reference ($\sim$4~minutes, needs Vitis HLS, no FPGA).

\subsection{Experiment customization}

Hardware mapping strategy (\texttt{no\_reuse}, \texttt{balanced}, \texttt{max\_reuse}),
tiling strategy, solver choice (\texttt{minlp}, \texttt{sa}, \texttt{ga},
\texttt{random}, \texttt{fixed}), DSP budget, objective weights and backend are all
fields of \texttt{STFlowConfig} (\texttt{stflow/config.py}), settable per test through
the category's \texttt{test\_config.yaml}. \texttt{STFLOW\_MODE} switches a run between
C-simulation, synthesis, co-simulation and bitstream execution. New workloads are added
as PyTorch modules under \texttt{test/}; the sweep grid is defined by \texttt{evgen}.

\subsection{Notes}

HLS synthesis, bitstream generation and on-board execution can all be re-run from
source, but they require a licensed Vitis HLS for synthesis, Vivado plus a physical Alveo U55C for the
board.

For MINLP-driven designs, a Tier~3 re-run on the same tool versions matches the archived
numbers exactly: regenerating \texttt{01\_mlp\_mixer} and \texttt{06\_random\_mlp\_seq}
end to end (all three mapping strategies, at the paper's 500-DSP budget) through DSE,
HLS synthesis and RTL co-simulation reproduced every predicted-latency, RTL-cycle and DSP
figure in \texttt{main.csv} cycle-for-cycle. What varies run to run: the SA and GA
searches are stochastic by construction (the GPT-2 study is 20 independent runs per
method), and Vitis HLS/Vivado are not bit-reproducible across tool patch levels, so on a
different patch level cycle counts, resource totals, achieved frequency and even which
designs route can move. The archived outputs are the reference for every reported
number; a reviewer with neither the tools nor the card loses no claim, since the full
implementation reports and on-board logs for all 425 experiments are archived.

Two provenance notes may surprise a reviewer. Per-design DSE reports are named
\texttt{reports/phoenix.csv} because the project was called Phoenix when the experiments
were run and was renamed to ST-Flow afterwards; the reporting code accepts either name.
And the CODO baseline designs were built from the third-party CODO artifact and
augmented only with AXI-DMA interfaces, so they share ST-Flow's off-chip execution
model --- 14 successful designs under \texttt{designs/} and 9 failures under
\texttt{failed/} are archived with their full Vitis reports, and
\texttt{scripts/codo/collect\_codo\_reports.py} re-derives \texttt{codo.csv} from them,
so Figure~5 is reproducible even though the CODO toolchain itself is not included.

\subsection{Methodology}

Submission, reviewing and badging methodology:
\begin{itemize}
  \item \url{https://www.acm.org/publications/policies/artifact-review-and-badging-current}
  \item \url{https://cTuning.org/ae}
  \item \url{https://pact2026.github.io/ae/}
\end{itemize}
